\documentclass[UTF8,fontset=none,11pt,a4paper]{ctexart}
\usepackage[margin=2.5cm]{geometry}
\usepackage{amsmath,amssymb}
\usepackage[most]{tcolorbox}
\usepackage{tabularx}
\usepackage{booktabs}
\usepackage[table]{xcolor}
\usepackage{enumitem}
\usepackage{microtype}
\usepackage[hidelinks]{hyperref}
\usepackage{fontspec}
\setCJKmainfont{Noto Serif CJK SC}[
  BoldFont={Noto Serif CJK SC Bold},
  ItalicFont={Noto Serif CJK SC}
]
\setCJKsansfont{Noto Sans CJK SC}[
  BoldFont={Noto Sans CJK SC Bold}
]

\definecolor{deepgreen}{HTML}{1B5E20}
\definecolor{lightgreen}{HTML}{E8F5E9}

\pagestyle{plain}
\hyphenation{normalization}
\setlength{\parindent}{0pt}
\setlength{\parskip}{0.35em}
\setlist[itemize]{leftmargin=*, itemsep=1pt, topsep=2pt, parsep=0pt}

\newtcolorbox{answerbox}{
  colback=lightgreen,
  colframe=deepgreen,
  boxrule=0.6pt,
  arc=1mm,
  left=4mm,
  right=4mm,
  top=2.5mm,
  bottom=2.5mm
}

\begin{document}

\begin{tcolorbox}[
  colback=white,
  colframe=deepgreen,
  boxrule=0.8pt,
  arc=1mm,
  left=4mm,
  right=4mm,
  top=3mm,
  bottom=3mm
]
\textcolor{deepgreen}{\textbf{Problem (think\_about\_length\_normalization):
Think about length normalization (1 point)}}
\par\medskip
Before running any experiments, think about the difference between normalizing
each sequence by sequence length, versus normalizing all sequences by the same
constant. What are the pros and cons of each approach? Are there specific
settings or examples where one approach seems better?
\par\smallskip
\textbf{Deliverable:} A few sentences of discussion.
\end{tcolorbox}

{\large\bfseries\color{deepgreen} 两种估计器}\par\smallskip

设 \(B\) 为 prompt 数，\(G\) 为组大小，\(y_{ij}\) 是第 \(i\) 个问题的第 \(j\) 个
回答，\(L_{ij}=|y_{ij}|\) 是它的长度，\(A_{ij}\) 是组归一化后的 advantage。
忽略基线选择和 advantage 归一化（它们与本题正交），batch 策略梯度为
\begin{equation*}
\nabla_\theta J(\theta)
=
\frac{1}{BG}\sum_{i=1}^{B}\sum_{j=1}^{G}
A_{ij}\sum_{t=1}^{L_{ij}}
\nabla_\theta \log \pi_\theta
\bigl(y_{ij,t}\mid x_i,y_{ij,<t}\bigr).
\end{equation*}
两种归一化只改变每个 token 的权重：
\begin{equation}
\tag{28}
\hat g_{\mathrm{seq}}
=
\frac{1}{BG}\sum_{i=1}^{B}\sum_{j=1}^{G}
\frac{1}{L_{ij}}\sum_{t=1}^{L_{ij}}
A_{ij}\,\nabla_\theta \log \pi_\theta
\bigl(y_{ij,t}\mid x_i,y_{ij,<t}\bigr),
\end{equation}
\begin{equation}
\tag{33}
\hat g_{\mathrm{const}}
=
\frac{1}{Z}\sum_{i=1}^{B}\sum_{j=1}^{G}\sum_{t=1}^{L_{ij}}
A_{ij}\,\nabla_\theta \log \pi_\theta
\bigl(y_{ij,t}\mid x_i,y_{ij,<t}\bigr),
\qquad Z=BGL.
\end{equation}

常数归一化给每个 token 相同的权重 \(1/Z\)；序列长度归一化给长度为 \(L_{ij}\) 的
序列中的 token 权重 \(1/(BG\,L_{ij})\)，权重取决于该序列自己的长度。如果所有
序列长度相同，两种估计器只差一个全局标量，学习率吸收后完全等价；真正的区别只在
长度变化时出现。

\begin{answerbox}
\textcolor{deepgreen}{\textbf{两种方法的对比}}

\smallskip
{\small
\arrayrulecolor{deepgreen}
\renewcommand{\arraystretch}{1.2}
\begin{tabularx}{\linewidth}{@{}>{\raggedright\arraybackslash}p{0.2\linewidth}
                                   >{\raggedright\arraybackslash}X
                                   >{\raggedright\arraybackslash}X@{}}
\toprule
\rowcolor{lightgreen}
 & \textbf{按序列长度归一化} & \textbf{按常数归一化} \\
\midrule
每个 token 的权重 & \(1/(BG\,L_{ij})\)，依赖该序列长度 & \(1/Z\)，所有 token 相同 \\
每个序列的影响 & 相同，均为 \(1/(BG)\) & 与该序列的 token 数成正比 \\
偏差 & 有偏：\(1/L_{ij}\) 不在真实梯度中 & 无偏：只差全局常数 \(BG/Z\) \\
方差与稳定性 & 较低；长序列不会主导更新 & 较高；长序列可能主导更新 \\
更适合 & 长度相近，或长度与 reward 无关 & 长度与正确率相关（如 GSM8K CoT） \\
\bottomrule
\end{tabularx}
}
\end{answerbox}

\begin{answerbox}
\textcolor{deepgreen}{\textbf{优缺点}}

\textbf{序列长度归一化。} 优点是每个 response 等权，长度分布变化时 loss 和梯度
尺度稳定，\mbox{长序列}不会因为 token 多而主导 batch，方差通常更低；log 出来的 loss
也可以跨 batch 比较。缺点是 \(1/L_{ij}\) 不在真实策略梯度中，估计有偏，而且偏差
不会随样本数增加而消失：长而正确的回答被强化得更少，长而错误的回答被抑制得更少，
隐式地偏向更短的输出。per-token 有效学习率和固定 clip 阈值的行为也依赖长度；如果
正确率与长 CoT 相关，这种偏置会抑制模型写出长推理链。

\textbf{常数归一化。} 优点是与推导一致，每个 token 等权，估计器等于真实策略梯度
乘一个全局常数 \(BG/Z\)，没有逐序列的长度偏差；长度中性，固定 \(Z=BGL\) 时归一化
常数在 batch 间不变，学习率调参更有意义。Dr. GRPO 正是去掉长度归一化（以及 std
归一化）来恢复估计器的正确期望。缺点是梯度大小随总 token 数增长，长回答较多的
batch 会产生更大的更新；在重尾长度分布下，少数长序列可能主导梯度并增大方差，
通常需要梯度裁剪和更仔细的学习率/\(Z\) 选择。
\end{answerbox}

\begin{answerbox}
\textcolor{deepgreen}{\textbf{什么时候哪种更好？}}
\begin{itemize}
  \item \textbf{长度几乎相同。} 两者只差全局标量，常数归一化更简单且无偏。
  \item \textbf{长度变化大、reward 与长度无关。} 序列归一化方差更低、优化尺度
        更稳定；常数归一化无偏但更 noisy。
  \item \textbf{reward 与长度相关（GSM8K 的 CoT）。} 序列归一化会偏向更短的输出，
        可能抑制真正解决问题的长推理链；这是常数归一化（Dr. GRPO）更合适的
        主要场景。
  \item \textbf{重尾长输出或小 batch。} 常数归一化需要裁剪和合适的 \(Z\)；如果
        首要目标是稳定性，序列归一化更稳；如果追求理论正确性，常数归一化是
        正确的估计器。
  \item \textbf{作业中的变体。} 标准 GRPO 用序列归一化；GRPO\_constant、
        Dr. GRPO、RFT、MaxRL 用常数归一化，可以直接用实验检验长度偏置是否
        真的影响 GSM8K。
\end{itemize}
\end{answerbox}

\textcolor{deepgreen}{\textbf{结论。}} 序列长度归一化用“等权 + 降方差”换来了
长度相关的偏差；常数归一化用“无偏 + 长度中性”换来了更高的方差和对长度分布的
敏感性。对于 GSM8K 这种“更长的 CoT 往往更容易答对”的任务，长度偏置很可能有害，
这正是 Dr. GRPO 采用常数归一化的动机；如果各回答长度相近、长度与 reward 无关，
两者几乎等价。

\end{document}